\chapter*{第 8 章\quad 黏性流体动力学基础 课后习题}
\addcontentsline{toc}{chapter}{第 8 章 黏性流体动力学基础 课后习题}
\markboth{818 课后习题}{第 8 章 黏性流体动力学基础}

\begin{tishi}
\textbf{题量与出处：}本章课后习题取自赵琴《流体力学与流体机械》第 8 章（书页 p183--p184），
共 16 题：选择题 1 题（习题 8.1，含 3 个小题）、证明题 1 题（习题 8.5）、
计算题 14 题（习题 8.2--8.4、8.6--8.16）。

\textbf{考情定位：}《8-1 考情分析》slide33 明确"第八章：黏性流体动力学基础
会涉及选择题和填空题考点"；《9-1 考点分析》亦称"第六到第八章只有少量填空/判断/选择，
掌握基本概念＋公式套用即可"。但本章的平板边界层摩擦阻力计算在历年真题计算题中三度出现
（2015 年计算第 9 题、2016 年计算第 10 题、2018 年计算第 9 题），且习题 8.11 与 2015 年
计算第 9 题、习题 8.12 与 2018 年计算第 9 题、习题 8.8 与 2022 年填空题、习题 8.16 与
2019 年计算第 7 题\textbf{题设数值完全相同}，故本章课后习题必须逐题过关。

\textbf{OCR 校订说明（各题题设数值一律未作改动，仅订正字形错读）：}
习题 8.4 题给的速度分布公式在原扫描件中整行不可辨，本册据紧随其后的习题 8.5 所证明的
三次多项式 $u/U_\infty=(3/2)(y/\delta)-(1/2)(y/\delta)^3$ 补入，并在解答册中给出
"三次多项式"与"抛物线 $u/U_\infty=2(y/\delta)-(y/\delta)^2$"两种口径；
习题 8.8 油的运动黏度指数丢失，据 2020 年真题同一题的影印校订为
$\nu_2=8\times10^{-5}\ \mathrm{m^2/s}$；习题 8.10 水的运动黏度原 OCR 作 $10^{-2}$，
按水的物性及同章各题订为 $\nu=10^{-6}\ \mathrm{m^2/s}$；习题 8.15 的密度单位
原 OCR 误作 $\mathrm{kg/m^2}$，订为 $\mathrm{kg/m^3}$。

\textbf{本章公式体系（教材 8.1--8.9 节）：}两平行平板间层流的精确解（纯剪切流、泊肃叶流、
库特流）$\to$ 边界层概念与分离 $\to$ 动量积分方程 $\to$ 平板层流/湍流/混合边界层
（$\delta=5.0x/\sqrt{\Rey_x}$、$\delta=0.37x/\Rey_x^{1/5}$、
$C_D=1.328/\sqrt{\Rey_l}$、$C_D=0.074/\Rey_l^{1/5}-1700/\Rey_l$）
$\to$ 绕流阻力 $D=C_D(\rho U^2/2)A$ 与自由沉降速度。
\end{tishi}

\begin{ti}
\sloppy
\emergencystretch=2em

\item \textbf{选择题}
\begin{xiaowen}
\item 汽车高速行驶时所受到的阻力主要来自于（\quad）。
\begin{choices}
\item 汽车表面的摩擦阻力
\item 地面的摩擦阻力
\item 空气对头部的碰撞
\item 尾部的旋涡
\end{choices}
\item 边界层内的流动特点之一是（\quad）。
\begin{choices}
\item 黏性力比惯性力重要
\item 黏性力与惯性力量级相等
\item 压强变化可忽略
\item 流动速度比外部势流小
\end{choices}
\item 边界层的流动分离发生在（\quad）。
\begin{choices}
\item 物体后部
\item 零压梯度区
\item 逆压梯度区
\item 后驻点
\end{choices}
\end{xiaowen}
\quad\textbf{重要度 \xstarfull{4}}\quad\yiju{E4 2019 年 818 判断题（平板边界层内沿厚度方向压强不变）；E4 2020 年 818 简答题第 3 题（绕流阻力的组成与减小阻力的措施）；E1 slide33"第八章：黏性流体动力学基础 会涉及选择题和填空题考点"；E6 教材 8.2、8.7 节}\quad\nandu{易}

\item 两平行平板间的泊肃叶流动如习题 8.2 图所示，平板间距为 $2h$，板长为 $L$ 且 $h\ll L$。
黏性不可压缩流体在恒定压强差 $(p_1-p_2)$ 的作用下沿 $x$ 方向流动，求平板间的流体运动速度分布。
\tuyi{习题 8.2 图：上、下两块水平放置的平行平板，板长为 $L$、间距为 $2h$（$h\ll L$）；
取 $x$ 轴沿流动方向，$y$ 轴垂直于板面、原点取在两板中间的中线上，故上板位于 $y=+h$、
下板位于 $y=-h$；进口断面压强为 $p_1$、出口断面压强为 $p_2$，且 $p_1>p_2$，
流体在恒定压强差作用下沿 $x$ 方向作定常层流流动。}
\quad\textbf{重要度 \xstarfull{3}}\quad\yiju{E4 2019 年 818 选择题（给出两平板间的速度分布、求通过平板的单宽流量，与本题同属两平板间层流速度分布）；E6 教材 8.1 节两平行平板间层流的精确解（泊肃叶流）；E1 slide33 第八章为选择填空题考点}\quad\nandu{中}

\item 上题中，若下板固定不动，上板以速度 $U$ 沿流动方向运动，求平板间的流体运动速度分布。
（坐标、符号同习题 8.2 图）
\quad\textbf{重要度 \xstarfull{3}}\quad\yiju{E4 2019 年 818 选择题（两平板间层流的速度分布与单宽流量）；E6 教材 8.1 节（纯剪切流与泊肃叶流叠加即得库特流）；E1 slide33 第八章为选择填空题考点}\quad\nandu{中}

\item 一长 $1.2\ \mathrm{m}$、宽 $0.6\ \mathrm{m}$ 的平板，顺流放置于速度为 $0.8\ \mathrm{m/s}$ 的定常水流中，
设平板上边界层内的速度分布为 $u/U_\infty=(3/2)(y/\delta)-(1/2)(y/\delta)^3$
（原扫描件此处公式整行不可辨，本册据习题 8.5 补入），其中 $\delta$ 为边界层厚度，
$y$ 为至平板的垂直距离。试求：（1）边界层厚度的最大值；（2）作用在平板上的单面阻力。
\quad\textbf{重要度 \xstarfull{2}}\quad\yiju{E6 教材 8.5 节动量积分方程与三次多项式速度分布的推导（$\delta=4.64x/\sqrt{\Rey_x}$、$C_D=1.293/\sqrt{\Rey_l}$）；E4 2020 年填空、2021 年计算第 5 题虽考平板边界层厚度与阻力，但题给的是湍流/混合边界层口径且给定了 $\Rey$ 判据；E2"第六到第八章掌握基本概念＋公式套用即可"；故压至两星}\quad\nandu{难}

\item 设一平板顺流放置于速度为 $U_\infty$ 的均匀来流中，如已知平板上层流边界层内的速度分布
$u(y)$ 可用 $y$（$y$ 为该点至板面的距离）的 3 次多项式表示，
试证明这一速度分布可表示为
\[
\frac{u}{U_\infty}=\frac{3}{2}\left(\frac{y}{\delta}\right)-\frac{1}{2}\left(\frac{y}{\delta}\right)^3
\]
其中 $\delta$ 为边界层厚度。
\quad\textbf{重要度 \xstarfull{2}}\quad\yiju{E6 教材 8.5 节动量积分方程（三次多项式速度分布的四个边界条件即本题依据）；E1 slide33"第八章：黏性流体动力学基础 会涉及选择题和填空题考点"；E4 历年真题均为选择/填空/判断与计算题，未出现此类纯证明题，故压至两星}\quad\nandu{中}

\item 水以来流速度 $u=0.2\ \mathrm{m/s}$ 顺流绕过一块平板。已知水的运动黏度
$\nu=1.145\times10^{-6}\ \mathrm{m^2/s}$，试求距平板前缘 $5\ \mathrm{m}$ 处的边界层厚度。
\quad\textbf{重要度 \xstarfull{4}}\quad\yiju{E4 2020 年 818 填空题（平板末端边界层厚度，按层流口径 $\delta/x=5.0/\sqrt{\Rey_x}$）；E4 2021 年计算第 5 题（平板尾端边界层厚度与两面摩擦阻力）；E4 2022 年填空题（水绕平板末端边界层的流态判断）；同考点在真题中出现 3 次}\quad\nandu{中}

\item 一平板顺流放置于均匀来流中，若将平板的长度增加一倍，试问：平板所受的摩擦阻力将增加几倍？
（设平板边界层内的流动为层流）
\quad\textbf{重要度 \xstarfull{4}}\quad\yiju{E4 2015 年计算第 9 题（同一平板两种放置方向下 $F_1/F_2$，考的正是 $F\propto b\sqrt{l}$ 的比例关系）；E4 2016 年计算第 10 题（船体平板摩擦阻力）；E6 教材 8.6.1 节层流平板 $C_D=1.328/\sqrt{\Rey_l}$}\quad\nandu{中}

\item 流体以速度 $u=0.8\ \mathrm{m/s}$ 绕一块长 $L=2\ \mathrm{m}$ 的平板流动，如果流体分别是水
（$\nu_1=10^{-6}\ \mathrm{m^2/s}$）和油（$\nu_2=8\times10^{-5}\ \mathrm{m^2/s}$），试求平板末端的边界层厚度。
\quad\textbf{重要度 \xstarfull{5}}\quad\yiju{E4 2022 年 818 填空题与本题数值完全相同（水以 $0.8\ \mathrm{m/s}$ 绕长 $2\ \mathrm{m}$ 平板、$\nu=10^{-6}\ \mathrm{m^2/s}$，问末端边界层流态）；E4 2021 年计算第 5 题；E6 教材 8.6.1 与 8.6.2 节（层流与湍流边界层厚度公式）}\quad\nandu{中}

\item 边长为 $1\ \mathrm{m}$ 的正方形平板放在速度 $v=1\ \mathrm{m/s}$ 的水流中，求边界层的最大厚度及双面摩擦阻力，
分别按全板都是层流或者都是湍流两种情况进行计算，水的运动黏度 $\nu=10^{-6}\ \mathrm{m^2/s}$。
\quad\textbf{重要度 \xstarfull{4}}\quad\yiju{E4 2021 年计算第 5 题（求平板尾端边界层厚度及平板两面摩擦阻力）；E4 2018 年计算第 9 题（平板边界层摩擦阻力）；E6 教材 8.6.1、8.6.2 节（层流与湍流平板的 $\delta$ 与 $C_D$）}\quad\nandu{中}

\item 水渠底面是一块长 $L=30\ \mathrm{m}$、宽 $b=3\ \mathrm{m}$ 的平板，水流速度 $u_0=6\ \mathrm{m/s}$，
水的运动黏度 $\nu=10^{-6}\ \mathrm{m^2/s}$，试求：（1）平板前面 $x=3\ \mathrm{m}$ 一段板面的摩擦阻力；
（2）长 $L=30\ \mathrm{m}$ 的板面的摩擦阻力。
\quad\textbf{重要度 \xstarfull{4}}\quad\yiju{E4 2021 年计算第 5 题与 E4 2018 年计算第 9 题（平板边界层摩擦阻力，均须先判 $\Rey$ 是否超过转捩临界值、再选层流/湍流/混合公式）；E6 教材 8.6.3 节平板混合边界层}\quad\nandu{难}

\item 一块面积为 $2\ \mathrm{m}\times8\ \mathrm{m}$ 的矩形平板放在速度 $u=3\ \mathrm{m/s}$ 的水流中，
水的运动黏度 $\nu=10^{-6}\ \mathrm{m^2/s}$，平板放置的方法有两种：以长边顺着流速方向，摩擦阻力为 $F_1$；
以短边顺着流速方向，摩擦阻力为 $F_2$。试求比值 $F_1/F_2$。
\quad\textbf{重要度 \xstarfull{5}}\quad\yiju{E4 2015 年 818 计算第 9 题与本题数值完全相同（$2\ \mathrm{m}\times8\ \mathrm{m}$ 平板、$3\ \mathrm{m/s}$、$\nu=10^{-6}\ \mathrm{m^2/s}$、求 $F_1/F_2$）；E4 2016 年计算第 10 题；E6 教材 8.6.2 节湍流平板 $C_D=0.074/\Rey_l^{1/5}$}\quad\nandu{中}

\item 平底船的底面可视为宽 $b=10\ \mathrm{m}$、长 $L=50\ \mathrm{m}$ 的平板，船速 $v=4\ \mathrm{m/s}$，
水的运动黏度 $\nu=10^{-6}\ \mathrm{m^2/s}$，如果平板边界层转捩临界雷诺数
$\Rey_{xcr}=5\times10^5$，试求克服边界层阻力所需的功率。
\quad\textbf{重要度 \xstarfull{5}}\quad\yiju{E4 2018 年 818 计算第 9 题与本题数值完全相同（平底船底面 $b=10\ \mathrm{m}$、$L=50\ \mathrm{m}$、$v=4\ \mathrm{m/s}$、$\nu=10^{-6}\ \mathrm{m^2/s}$、$\Rey_{xcr}=5\times10^5$、求功率）；E4 2016 年计算第 10 题；E6 教材 8.6.3 节混合边界层}\quad\nandu{难}

\item 有 $45\ \mathrm{kN}$ 的重物从飞机上投下，要求落地速度不超过 $10\ \mathrm{m/s}$，重物挂在一张阻力系数
$C_D=2$ 的降落伞下面，不计伞重，设空气密度为 $\rho=1.2\ \mathrm{kg/m^3}$，求降落伞应有的直径。
\quad\textbf{重要度 \xstarfull{3}}\quad\yiju{E6 教材 8.7.2 节绕流阻力式 $D=C_D(\rho u^2/2)A$（$A$ 为迎流投影面积）；E4 2020 年简答题第 3 题（绕流阻力与减阻）；E1 slide33"第八章 会涉及选择题和填空题考点"；E2"掌握基本概念＋公式套用即可"}\quad\nandu{中}

\item 汽车以 $80\ \mathrm{km/h}$ 的时速行驶，其迎风面积为 $A=2\ \mathrm{m^2}$，阻力系数为 $C_D=0.4$，
空气的密度为 $\rho=1.25\ \mathrm{kg/m^3}$，试求汽车克服空气阻力所消耗的功率。
\quad\textbf{重要度 \xstarfull{3}}\quad\yiju{E6 教材 8.7.2 节（绕流阻力与克服阻力所需功率）；E4 2016 年计算第 10 题（船体阻力与功率，同型）；E1 slide33 第八章为选择填空题考点；E2"掌握基本概念＋公式套用即可"}\quad\nandu{易}

\item 炉膛的烟气以速度 $V_0=0.5\ \mathrm{m/s}$ 向上腾升，气体的密度为 $\rho=0.25\ \mathrm{kg/m^3}$，
动力黏性系数 $\mu=5\times10^{-5}\ \mathrm{N\cdot s/m^2}$，粉尘的密度 $\rho'=1200\ \mathrm{kg/m^3}$，
试估算此烟气能带走多大直径的粉尘？
\quad\textbf{重要度 \xstarfull{4}}\quad\yiju{E5 题库计算第 3 题（已知煤粉炉炉膛上升烟气最小速度 $0.5\ \mathrm{m/s}$、烟气 $\nu=230\times10^{-6}\ \mathrm{m^2/s}$，问 $d=0.1\ \mathrm{mm}$ 的煤粉是否沉降）与 E6 教材【例 8.4】为同源问题；E4 2019 年计算第 7 题（钢球在油中自由沉降）同属自由沉降考点；E6 教材 8.7.2 节式(8.63)}\quad\nandu{难}

\item 使小钢球在油中自由沉降以测定油的黏度。已知油的密度 $\rho=900\ \mathrm{kg/m^3}$，
小钢球的直径 $d=3\ \mathrm{mm}$，密度 $\rho'=7788\ \mathrm{kg/m^3}$，若测得钢球最终的沉降速度
$u=12\ \mathrm{cm/s}$，求油的动力黏度。
\quad\textbf{重要度 \xstarfull{5}}\quad\yiju{E4 2019 年 818 计算第 7 题与本题数值完全相同（$\rho=900\ \mathrm{kg/m^3}$、$d=3\ \mathrm{mm}$、$\rho'=7788\ \mathrm{kg/m^3}$、$v=12\ \mathrm{cm/s}$，求油的动力黏度，答案约 $0.281\ \mathrm{Pa\cdot s}$）；E6 教材 8.7.2 节圆球自由沉降的力平衡与 $C_D=24/\Rey$}\quad\nandu{中}

\end{ti}
